\section{Complete native enumeration}
\label{sec:complete-enumeration}

The lift construction of Section~\ref{sec:quotient-lift} produces a
subgroup
\[
L\leq\operatorname{Aut}(G60)
\]
of order $480$. To prove that $L$ is the full automorphism group, we
perform a second calculation that does not assume preservation of the
$a$-fiber partition.

\subsection{Independent graph-automorphism search}

The native graph is reconstructed directly from its exact sixty-vertex
edge list. The resulting graph has
\[
60\text{ vertices},
\qquad
120\text{ edges},
\qquad
\deg(v)=4
\]
for every vertex $v$, and is connected.

A complete graph-isomorphism enumeration is then run from the native
graph to itself. Each returned isomorphism is converted into a
permutation of the fixed native vertex ordering. Duplicate permutations
are removed by exact tuple comparison.

The enumeration returns
\[
480
\]
distinct automorphisms.

\begin{theorem}
\label{thm:native-order}
The native graph satisfies
\[
\left|\operatorname{Aut}(G60)\right|=480.
\]
\end{theorem}

\begin{proof}
The exhaustive native enumeration returns $480$ distinct permutations
preserving the complete edge set. Hence
\[
\left|\operatorname{Aut}(G60)\right|\leq480.
\]
Section~\ref{sec:quotient-lift} constructs a subgroup
\[
L\leq\operatorname{Aut}(G60)
\]
with $|L|=480$. Therefore
\[
\left|\operatorname{Aut}(G60)\right|=480.
\]
\end{proof}

\subsection{Equality with the lifted group}

Both the native enumeration and the quotient-lift construction produce
explicit permutations in the same sixty-vertex labeling. Their element
sets can therefore be compared directly.

\begin{theorem}
\label{thm:lifted-equals-full}
The lifted group is the full automorphism group:
\[
L=\operatorname{Aut}(G60).
\]
\end{theorem}

\begin{proof}
Every element of $L$ is a native automorphism, so
\[
L\subseteq\operatorname{Aut}(G60).
\]
Both groups have order $480$ by Theorems~\ref{thm:two-lifts} and
\ref{thm:native-order}. Hence equality follows.

Equivalently, direct comparison of the two explicit permutation sets
gives exact equality, with no element occurring in only one set.
\end{proof}

Thus every native automorphism arises as a lift of a quotient
automorphism.

\subsection{Preservation of the canonical fibers}

The native enumeration also permits a direct test of the quotient
structure. For every
\[
g\in\operatorname{Aut}(G60)
\]
and every native vertex $v$, the pair
\[
\{g(v),g(a(v))\}
\]
is again an $a$-fiber.

Consequently, every native automorphism induces a permutation of the
thirty quotient fibers.

\begin{proposition}
\label{prop:fiber-preservation}
Every automorphism of $G60$ preserves the $a$-fiber partition.
\end{proposition}

This proposition is stronger than the cardinality argument alone. It
shows that the quotient used to construct the lifted subgroup is not
merely convenient: it is respected by the full native automorphism
group.

\subsection{Centrality of the deck involution}

Fiber preservation does not automatically imply that an automorphism
commutes with the deck involution. This is checked directly on the
native permutations.

For every
\[
g\in\operatorname{Aut}(G60),
\]
one has
\[
ga=ag.
\]

\begin{proposition}
\label{prop:deck-central}
The canonical deck involution lies in the center:
\[
a\in Z(\operatorname{Aut}(G60)).
\]
\end{proposition}

The quotient action therefore gives an exact sequence
\[
1
\longrightarrow
\langle a\rangle
\longrightarrow
\operatorname{Aut}(G60)
\longrightarrow
\operatorname{Aut}(G30)
\longrightarrow
1,
\]
with central kernel of order $2$.

\subsection{What the enumeration proves}

The independent calculation closes three logically distinct questions.

First, there are no native automorphisms outside the $480$ lifted
permutations.

Second, the canonical $a$-fiber system is invariant under the full
automorphism group.

Third, the deck involution is central in the full group.

The remaining problem is no longer to find additional automorphisms. It
is to identify the abstract structure of the already complete
$480$-element group. We turn to that problem in the next section.
